% =====================================================================
% Modelo ABNT de relatorio tecnico e/ou cientifico
% EguaTex (www.eguatex.com.br). Licenca CC0 1.0: dominio publico.
%
% Segue a ABNT NBR 10719:2015 e as normas que ela chama: NBR 6023:2025
% (referencias), NBR 6024:2012 (secoes), NBR 6027:2012 (sumario),
% NBR 6028:2021 (resumo) e NBR 10520:2023 (citacoes).
%
% Como usar: troque os textos de exemplo pelo seu conteudo e compile com
% pdfLaTeX e BibTeX (o EguaTex faz isso sozinho, e as obras citadas ficam
% no arquivo referencias.bib). A formatacao fica no arquivo
% abnt-eguatex.sty. O relatorio usa espacamento simples (5.1).
% =====================================================================
\documentclass[12pt,a4paper,oneside]{article}
\usepackage[relatorio]{abnt-eguatex}

\begin{document}

% =====================================================================
% CAPA (opcional, 4.1.1.1)
% Recomenda-se: instituicao responsavel e endereco, numero do relatorio,
% ISSN (se houver), titulo e subtitulo e classificacao de seguranca (se
% houver). Apague este bloco se nao quiser capa.
% =====================================================================
\begin{titlepage}\centering\setstretch{1}
  NOME DA INSTITUIÇÃO RESPONSÁVEL\par
  Endereço completo da instituição\par
  \vspace{2\baselineskip}
  Relatório nº 00/2026\par
  ISSN 0000-0000\par
  \vfill
  {\bfseries TÍTULO DO RELATÓRIO:} subtítulo, se houver\par
  \vfill
  Classificação de segurança: ostensivo\par
\end{titlepage}
\setcounter{page}{1}

% =====================================================================
% FOLHA DE ROSTO, ANVERSO (obrigatoria, 4.2.1.1.1)
% Ordem: entidade responsavel, titulo do projeto ou programa, titulo do
% relatorio, subtitulo, volume, codigo de identificacao, classificacao de
% seguranca, autor, local e ano.
% =====================================================================
\begingroup\thispagestyle{empty}\centering
  NOME DO ÓRGÃO OU ENTIDADE RESPONSÁVEL\par
  \vspace{2\baselineskip}
  Projeto: Nome do projeto, programa ou plano\par
  \vfill
  {\bfseries TÍTULO DO RELATÓRIO:} subtítulo, se houver\par
  \vspace{\baselineskip}
  Volume 1\par
  \vspace{\baselineskip}
  Código de identificação: SIGLA-RT-2026-001\par
  Classificação de segurança: ostensivo\par
  \vfill
  Nome do autor\par
  Titulação ou função do autor\par
  \vfill
  Cidade\par
  2026\par
\endgroup
\clearpage

% =====================================================================
% FOLHA DE ROSTO, VERSO (4.2.1.1.2)
% Equipe tecnica (opcional) e ficha catalografica (opcional, mas
% obrigatoria quando nao houver formulario de identificacao no final).
% =====================================================================
\thispagestyle{empty}
\noindent\textbf{Equipe técnica}\par\vspace{0.5\baselineskip}
\noindent Coordenação: Nome, titulação\\
Pesquisa: Nome, titulação\\
Colaboração: Nome, titulação
\vspace*{\fill}
\begin{center}
  \fbox{\begin{minipage}[c][7.5cm][c]{12.5cm}\centering\abntmenor\setstretch{1}
    Espaço da ficha catalográfica.\\[4pt]
    Peça a ficha à biblioteca da instituição. Ela é feita por bibliotecário\\com registro no CRB.
  \end{minipage}}
\end{center}
\clearpage\addtocounter{page}{-1}

% Errata (opcional, 4.2.1.2): so depois do relatorio impresso.

% Agradecimentos (opcional, 4.2.1.3)
\begin{agradecimentos}
  Aqui vão os agradecimentos a quem contribuiu de forma relevante.
\end{agradecimentos}

% Resumo (obrigatorio, 4.2.1.4; NBR 6028:2021): 150 a 500 palavras
\begin{resumo}
Aqui vai o resumo do relatório, em um único parágrafo, com frases curtas e o verbo na terceira pessoa. Ele informa o objetivo, o método, os resultados e as conclusões. Em relatórios técnicos e científicos, o resumo tem de 150 a 500 palavras.
\palavraschave{palavra-chave; segunda palavra-chave; terceira palavra-chave.}
\end{resumo}

% Listas (opcionais, 4.2.1.5 a 4.2.1.8). Apague as que nao usar.
\listadefiguras
\listadetabelas

% Sumario (obrigatorio, 4.2.1.9)
\sumario

% =====================================================================
% ELEMENTOS TEXTUAIS (4.2.2). O nome das secoes fica a seu criterio.
% =====================================================================
\textual

\section{Introdução}

A introdução apresenta os objetivos do relatório e as razões de sua elaboração. O relatório usa letra tamanho 12 e espaçamento simples.

\section{Desenvolvimento}

O desenvolvimento detalha a pesquisa ou o estudo realizado.

\subsection{Materiais e métodos}

Aqui vão os materiais e os métodos. As citações seguem o sistema autor-data, como em \cite[p.~15]{livro2024} ou \citeonline{pagina2025}. No código, a citação é feita com \texttt{\textbackslash cite\{chave\}}, e a obra entra sozinha na lista de referências.

\subsection{Resultados}

Toda ilustração é citada no texto, como a Figura~\ref{fig:exemplo}, com o título em cima e a fonte embaixo.

\begin{figure}[htbp]
  \begin{ilustracao}{0.7\textwidth}
    \caption{Título da figura}\label{fig:exemplo}
    \fbox{\parbox[c][4cm][c]{\dimexpr\linewidth-2\fboxsep-2\fboxrule}{\centering Coloque a imagem aqui com\\\texttt{\textbackslash includegraphics[width=\textbackslash linewidth]\{arquivo\}}}}
    \fonte{elaborada pelo autor (2026).}
  \end{ilustracao}
\end{figure}

As tabelas seguem as normas de apresentação tabular do IBGE, como a Tabela~\ref{tab:exemplo}.

\begin{table}[htbp]
  \begin{ilustracao}{0.7\textwidth}
    \caption{Título da tabela, com o que, onde e quando}\label{tab:exemplo}
    \centering
    \begin{tabular*}{\linewidth}{@{\extracolsep{\fill}}lrr}
      \toprule
      Amostra & Medida & Variação (\%) \\
      \midrule
      Amostra 1 & 12,5 & 2,1 \\
      Amostra 2 & 13,1 & 1,8 \\
      \bottomrule
    \end{tabular*}
    \fonte{elaborada pelo autor (2026).}
  \end{ilustracao}
\end{table}

\section{Considerações finais}

As considerações finais apresentam as conclusões do relatório.

% =====================================================================
% ELEMENTOS POS-TEXTUAIS (4.2.3)
% =====================================================================
% Referencias (obrigatorias quando ha citacoes, NBR 10719:2015, 4.2.3.1;
% formato da NBR 6023:2025). A lista e montada sozinha
% a partir do arquivo referencias.bib: entram so as obras citadas no texto,
% em ordem alfabetica e no formato da norma.
\bibliography{referencias}

\apendice{Título do apêndice}
Aqui vai o conteúdo do apêndice, como planilhas de dados. Apague se não houver.

% Formulario de identificacao (opcional, 4.2.3.6). Obrigatorio quando o
% relatorio nao tiver ficha catalografica.
\titulopostextual{Formulário de identificação}
{\setstretch{1}\abntmenor
\noindent\begin{tabular}{|p{0.47\textwidth}|p{0.47\textwidth}|}
  \hline
  \multicolumn{2}{|c|}{\textbf{Dados do relatório técnico e/ou científico}} \\ \hline
  Título e subtítulo: & Classificação de segurança: \\ \hline
  Tipo de relatório: & Data: \\ \hline
  Título do projeto, programa ou plano: & Nº: \\ \hline
  \multicolumn{2}{|p{0.96\textwidth}|}{Autor(es):} \\ \hline
  \multicolumn{2}{|p{0.96\textwidth}|}{Instituição executora e endereço completo:} \\ \hline
  \multicolumn{2}{|p{0.96\textwidth}|}{Instituição patrocinadora e endereço completo:} \\ \hline
  \multicolumn{2}{|p{0.96\textwidth}|}{Resumo:\vspace{2cm}} \\ \hline
  \multicolumn{2}{|p{0.96\textwidth}|}{Palavras-chave ou descritores:} \\ \hline
  Edição: & Nº de páginas: \\ \hline
  Nº do volume ou parte: & Nº de classificação: \\ \hline
  ISSN: & Tiragem: \\ \hline
  Preço: & Distribuidor: \\ \hline
  \multicolumn{2}{|p{0.96\textwidth}|}{Observações ou notas:} \\ \hline
\end{tabular}\par}

\end{document}
